\documentclass[11pt]{article}
\usepackage[utf8]{inputenc}
\usepackage[T1]{fontenc}
\usepackage{lmodern}
\usepackage[margin=1in]{geometry}
\usepackage{csquotes}
\usepackage[style=authoryear,backend=biber,sorting=nyt,maxbibnames=99,giveninits=true,uniquename=init]{biblatex}
\usepackage[hidelinks]{hyperref}
\addbibresource{references.bib}
\begin{document}
\section*{Phosphorus Loss Pathways and Soil Controls}

\subsection{Overview}

Phosphorus (P) leaves agricultural fields along two hydrologically distinct
pathways---surface runoff and subsurface (tile) drainage---and the environmental
significance of each depends on how P is partitioned between them, on the balance
of dissolved and particulate fractions within each pathway, and on the soil
properties that govern P release and retention. The evidence assembled here
establishes several consolidated conclusions. First, subsurface drainage is not
the ``clean'' pathway it was once assumed to be: across humid, cold-temperate,
and tile-intensive landscapes it can dominate edge-of-field water export and
carry a large, sometimes majority, share of both dissolved reactive P (DRP) and
total P (TP) loads \parencite{Pease2018,VanEsbroeck2016,Qi2018,Sadhukhan2019b}.
Second, particulate P (PP) frequently dominates drainage P losses, a pattern that
overturned the older expectation of purely dissolved subsurface transport and is
tied mechanistically to preferential (macropore) flow of fine topsoil particles
to drains \parencite{Uusitalo2001,Lu2003,Laubel1999,Sims1998}. Third, dissolved P
concentrations and losses are strongly and predictably controlled by soil P
status---soil test P (STP), degree of P saturation (DPS), and P-sorption
capacity---so that soils enriched by long-term fertilisation or manuring, together
with sandy and organic soils of low sorption capacity, present the greatest
subsurface P-loss risk \parencite{Sharpley1995,Beauchemin2003,Siemens2004,Sims1998,McDowell2001b}.
Fourth, P export is highly episodic and seasonally concentrated: a small fraction
of the year, dominated by storm events, snowmelt, and the non-growing season,
carries the great majority of annual load
\parencite{King2015,VanEsbroeck2016,Brendel2019,Grant1996}. These lines of
evidence, spanning plot lysimeters, edge-of-field monitoring, watershed sediment
budgets, and process models, are developed in turn below, along with what is known
about soil P pools and speciation, organic and redox-active soils, management and
fertilisation controls, P-index and modelling frameworks, and edge-of-field
retention.

\subsection{Partitioning of Phosphorus Losses Between Surface Runoff and
Subsurface Drainage}

\subsubsection{Water partitioning and the hydrologic dominance of tile flow}

Tile drainage commonly captures the largest share of field water export. In
edge-of-field studies in the Eastern Corn Belt, subsurface discharge accounted for
$81 \pm 23\%$ of annual discharge \parencite{Pease2018}, and in humid
cold-temperate Ontario tile drains contributed 78--90\% of annual water export
with surface runoff supplying only 10--22\% \parencite{VanEsbroeck2016}. Installing
or improving subsurface drainage fundamentally re-routes water: improved drainage
in Finland shifted runoff from being surface-dominated (drainage 10--40\% of total
runoff before) to drainage-dominated (50--90\% after), with more than 80\% of
precipitation water evacuated by tiles and less than 10\% by surface runoff on
ploughed clay immediately after installation \parencite{Turtola1995}. Consistent
with this, subsurface drainage reduced surface runoff by an average of 35\% over a
ten-year study on Lower Mississippi Valley alluvial soil \parencite{BENGTSON1995},
and tile-drained plots in New York generated much greater cumulative runoff overall
but lower \emph{surface} runoff because most water left via tiles
\parencite{Klaiber2020}. Growing-season tile drainage generally tracked surface
water discharge but reached 80\% of surface water volume in a Massachusetts
cranberry system \parencite{Millar2025}, and in an intensively tiled Illinois
watershed subsurface drainage and surface runoff removed averages of 16.1\% and
2.6\% of rainfall, respectively \parencite{Algoazany2007}.

\subsubsection{Concentrations, loads, and pathway partitioning of P}

A recurring and important distinction is that surface runoff usually carries higher
P \emph{concentrations}, whereas tile drainage---because it carries more
water---can dominate annual \emph{loads}. Flow-weighted DRP and TP concentrations
were significantly greater in surface runoff than in subsurface discharge in Ohio,
yet annual DRP and TP loads were significantly greater from subsurface discharge,
which supplied 71\% of DRP and 69\% of TP load \parencite{Pease2018}. The same
concentration--load contrast appears in long-term Swedish monitoring, where TP
concentrations exceeded those of drainage in all but 3 of 32 years (32-year means
0.27 vs. 0.04~mg~L$^{-1}$) and annual P losses were greater via surface runoff,
while N losses were greater via drainage \parencite{Norberg2022}. In Ontario, tiles
exported 40--77\% of TP load and 19--67\% of DRP load
\parencite{VanEsbroeck2016}; in Wisconsin, tiles contributed 21--52\% of TP and
21--68\% of DRP field losses (concluded as 17--41\% of total P loss)
\parencite{Madison2014}; and a southern-Ontario field lost 65.4\% of annual TP
(2.85 of 4.36~kg~P~ha$^{-1}$) via tiles \parencite{Qi2018}. Regional syntheses put
subsurface drainage at 12--60\% of TP transfer versus 40--88\% for surface erosion
in north-west Europe \parencite{Ulén2007}, and a text-mining review reported tiles
transporting up to 77\% of TP, 67\% of soluble reactive P, and 63\% of watershed
suspended sediment \parencite{Elsayed2023}. Models corroborate the field
partitioning: process simulations attributed 66\% of DRP and 76\% of TP loss to
subsurface drainage against observed values of 70\% and 80\% \parencite{Grewal2025},
and RZWQM2-P reproduced tile drainage as the major DRP pathway with PP split roughly
equally between runoff and drainage \parencite{Sadhukhan2019b,Sadhukhan2019}.

\subsubsection{Settings where surface runoff dominates}

The dominance of tiles is not universal, and the counter-examples define the
conditions under which surface runoff prevails. In snowmelt-affected clay
(vertisolic) landscapes of the northern Great Plains, overland flow contributed
72--89\% of runoff and $>95\%$ of annual SRP and TP losses, because tiles were
decoupled from the surface by frozen ground and carried water only in wetter
periods; SRP and TP were 600\% and 500\% greater in overland flow than tile flow
\parencite{Kokulan2019b,Kokulan2021}. Surface runoff was also the major pathway for
sediment, TP, and soluble P (though not nitrate) under Minnesota tillage systems
\parencite{Zhao2001}, was the dominant P pathway from perennial pastures in Victoria
when runoff occurred \parencite{Melland2008}, and modelling of an Ohio field
attributed 78\% of predicted TP to surface runoff \parencite{Askar2021b}. In flooded
paddy systems, surface runoff exceeded subsurface pathways for P during rice growth,
and shallow subsurface drainage lowered TP loss mainly by reducing flooding
\parencite{DarziNaftchali2016}. Thus pathway partitioning is strongly modulated by
climate (snowmelt vs.\ rain-driven), soil texture, and drainage design rather than
being a fixed property.

\subsection{Dissolved and Particulate Phosphorus Fractions and Bioavailability}

\subsubsection{Particulate P as a major, often dominant, drainage fraction}

A central state-of-knowledge finding is that PP is frequently the dominant P form in
tile drainage, contrary to the earlier presumption that filtration by the soil
profile leaves subsurface water dissolved-dominated. In Finnish clay soils, PP was
the grand-mean 92\% of TP in both surface runoff and drainage, and most PP (93\%)
was nondesorbable \parencite{Uusitalo2001}; across sites, PP comprised 73--94\% of
runoff P \parencite{Uusitalo2003}. At Rothamsted's Broadbalk drains, PP averaged 54\%
of TP (vs.\ 36\% for molybdate-reactive P) and reached 81.1\% on individual plots,
with PP strongly linearly related to particulate matter (PM), which reached
$\sim$0.6~g~L$^{-1}$ \parencite{Lu2003,Lu2004}. PP dominated TP export in both tile
drainage and surface runoff in Ontario (though DRP:TP was higher in surface runoff)
\parencite{VanEsbroeck2016}, made up most of drainage TP in Massachusetts
\parencite{Millar2025}, and dominated Danish catchment losses (55--71\% of TP at
three of four catchments) \parencite{Grant1996}. Catchment monitoring likewise showed
drainage TP dominated by PP (only 39.8\% as DRP) with PP losses exceeding dissolved
losses \parencite{Brendel2019}, and level Dutch grassland lost nutrients dominated by
particulate forms via trenches and pipe drains \parencite{vanderSalm2012}. Over 80\%
of predicted surface and subsurface P losses were particulate in an Ohio simulation
\parencite{Askar2021b}, and pasture P export in New Zealand was $\sim$80\%
particulate-associated \parencite{Gillingham2000}. Even where TP was small, PP
dominated: a Finnish controlled-drainage peat field lost only
0.30~kg~P~ha$^{-1}$~yr$^{-1}$ of TP but 71\% as PP \parencite{Pham2023}.

\subsubsection{Dissolved-dominated systems and DRP importance}

Dissolved forms nonetheless dominate in specific settings, particularly
snowmelt-driven, grassland, high-STP, and organic-soil systems. In north-west
Europe the dissolved fraction spanned 9--93\% of TP \parencite{Ulén2007}. Tile SRP
averaged 60\% of TP (median 69\%) across Michigan sites \parencite{Clement2017}, and
DRP dominated fertilised continuous-grass tile losses (72\% of TP) even as PP
dominated corn and unfertilised systems \parencite{Zhang2015b}. In Scotland,
molybdate-reactive P accounted for only 42\% of subsurface TP loss, with non-MRP
dissolved forms 41\% and PP $\sim$17\%, underscoring the need to measure all P forms
\parencite{Hooda1999}. Organic soils in New Zealand and Ireland produced very high
filterable P loads and high overland-flow DRP concentrations, respectively
\parencite{McDowell2015b,Kurz2005}, and organic soils generally carry less
particulate but potentially higher TP concentrations than mineral soils
\parencite{Grenon2021}. Watershed orthophosphate:TP ratios rose in tile-drained and
baseflow-dominated Iowa systems \parencite{Schilling2017b}. Dissolved organic P
(DOP) is a consistently smaller but non-trivial fraction: DOP was 0.01--0.1~mg~L$^{-1}$
and generally below MRP at Rothamsted \parencite{Lu2003,Lu2004}, whereas the non-MRP
(DOP) fraction reached 44\% of total dissolved P in Scottish grassland
\parencite{Hooda1999}, and DOP leaching increased with peat thickness in Finland
\parencite{YliHalla2022}.

\subsubsection{Bioavailability of transported P}

Because algal availability, not merely total mass, determines eutrophication risk,
several studies quantified bioavailable P. Desorbable PP (PPi) and DRP were of equal
importance to potentially algal-available P in Finnish clays, with PPi accounting for
47\% of anion-exchange-resin P \parencite{Uusitalo2001}. Bioavailable estimates
depended strongly on the extraction: DRP + AER-PP yielded 62--669~g~ha$^{-1}$ across
sites, whereas the more aggressive DRP + BD-PP gave 145--1560~g~ha$^{-1}$, and PP
contributed significantly to bioavailable loss in both pathways, highlighting erosion
control as a water-protection measure \parencite{Uusitalo2003}. Only 8--18\% of TP
export from grassland lysimeters was molybdate-reactive \parencite{Bilotta2008}, and
in New Zealand streams bioavailable and total P in trapped sediment were highest at
downstream sites and were enriched when little sediment was trapped
\parencite{McDowell2004}. Sediment P enrichment declined with increasing sediment
yield (from $\sim$2000--4000 to $\sim$300--1000~mg~P~kg$^{-1}$), so total P yield was
not closely tied to sediment yield \parencite{Gillingham2000}.

\subsection{Preferential Flow, Sediment Mobilisation, and Particulate P Transport}

\subsubsection{Macropore flow as the dominant particulate-transport mechanism}

Preferential flow through cracks, fissures, macropores, and worm channels is the
mechanism reconciling substantial subsurface PP loss with the filtering expected of
matrix flow. It is identified as a major, rapid pathway delivering water and P to
drains \parencite{King2015,Sims1998}, and $^{137}$Cs tracing repeatedly shows that
drainage PM originates from topsoil: median grain sizes of 2--5~$\mu$m indicate fine
topsoil particles translocated to drains within hours of rainfall, creating a direct
link between topsoil and surface waters \parencite{Laubel1999}. Surface-runoff and
subsurface sediments were nearly identical in PPi content and $^{137}$Cs activity,
confirming macropore transport of surface soil particles through the profile, with
enrichment ratios of $\sim$2 for fine particles \parencite{Uusitalo2001}. Substantial
P losses occurred in Scottish grassland despite high subsoil P-sorption capacity,
implying preferential flow bypass \parencite{Hooda1999}, and macropore flow through
shrinkage cracks likely explained much DRP loss through Dutch pipe drains
\parencite{vanderSalm2012}. In vertisolic soils, preferential flow during dry summer
and fall raised tile P concentrations, though low flow limited loads
\parencite{Kokulan2021}. Ridge tillage increased subsurface sediment loss relative to
moldboard plowing, consistent with preferential flow \parencite{Zhao2001}. Field-scale
mapping links preferential transport to physical soil properties: in a Danish field
the clay-rich, high-bulk-density, low-organic-carbon northern third showed fast 5\%
tritium arrival, high tracer recovery (mean 67\%), and elevated particle and PP
leaching, identifying it as a leaching-vulnerable zone \parencite{Norgaard2013}.
First-flush particle concentrations of 62--390~mg~L$^{-1}$ carried PP of
0.07--1.73~mg~L$^{-1}$ before stabilising, with accumulated particle mass correlated
($r=0.66$) to first-flush PP \parencite{Norgaard2013}. The proportion of macropore
flow (MF) is recognised as the greatest uncertainty in representing tile P export in
indices, warranting further quantification across soils \parencite{Reid2012}.

\subsubsection{Event-driven particulate P and drainage as a sediment source}

PP loss is episodic and event-driven. PM, PP, and DP concentrations peaked in initial
drainage flow, and nearly one-third of annual PM and PP loss occurred during rapid
drainage in sampled storms \parencite{Laubel1999}. Storm drainage at a Danish
catchment was strongly enriched in DP and PP, PM peaked before discharge (165~mg~L$^{-1}$
$\approx$ 0.69~mg~P~L$^{-1}$ PP) suggesting exhaustion of fine-particle supply, and PP
was linearly correlated with PM at all catchments \parencite{Grant1996}. Because of
this episodicity, discrete sampling underestimated annual P loss by 26\% to $>$100\%
relative to intensive storm sampling, so reliable estimates require event-based
strategies \parencite{Grant1996}. Turbidity served as a strong suspended-solids proxy
in Swedish drainage \parencite{Ulén2018}, and grassland lysimeters exported up to
14.85~kg suspended solids and 50~g TP per event, with drained land exporting up to
52\% less than undrained land \parencite{Bilotta2008}.

\subsubsection{Catchment sediment budgets and tile drains as sediment sources}

At catchment scale, tile drains are quantitatively major sediment pathways.
Source-fingerprinting attributed 28--55\% of catchment sediment yield to field drains
across UK subcatchments, with eroding channel banks a smaller source (6--12\%) and low
overall sediment delivery ratios (14--27\%) indicating substantial in-catchment
redeposition \parencite{Walling2002}. In monitored UK drains, topsoil supplied more
than two-thirds of the sediment load (70.9--88.6\%), sediment load correlated with
drainage volume rather than rainfall, and drains contributed over 50\% of the annual
sediment budget in one catchment \parencite{Chapman2005}. Sediment yields ranged
70--978~kg~ha$^{-1}$~yr$^{-1}$ with PP yields 0.04--1.1~kg~ha$^{-1}$~yr$^{-1}$
\parencite{Chapman2005}. Tile-drain sediment flocculates under shear, with maximum
flocculation at 0.169~N~m$^{-2}$ and a critical deposition shear stress of
0.056~N~m$^{-2}$, information relevant to its fate in receiving streams
\parencite{Stone1997}. Clayey soils exported measurable sediment-associated organic
carbon and N (up to 35\% and 20\% of loads) chiefly via subsurface drainage, whereas a
sandy soil eroded negligibly and exported only aqueous forms; discharge and soil
disturbance (recent plowing) controlled export \parencite{Manninen2023}. At larger
scales, non-field sources can dominate: streambanks and ravines supplied 76\% of
sediment at a Minnesota outlet \parencite{Dalzell2018}, and streambed storage exceeded
annual streambank erosion and suspended-sediment load in an Ohio lake-plain tributary,
where streambanks were the largest streambed-sediment source and much sediment had long
residence times \parencite{Williamson2025}.

\subsection{Soil Controls: Soil Test P, Sorption Capacity, and Degree of P Saturation}

\subsubsection{Soil test P and dissolved P release}

Dissolved P loss is tightly coupled to soil P status. Loss of P in surface runoff and
subsurface drainage is strongly related to soil P content
($r^2 = 0.85$--0.96 for surface water, 0.86 for subsurface), and at tile-drain depth
(40--50~cm) NaCl-extractable soil P related to storm drainage loss ($r^2=0.86$)
\parencite{Sharpley1995}. Higher STP was associated with higher DRP in both tiles and
runoff, while soil texture had limited overall influence in a large Ohio edge-of-field
dataset \parencite{Osterholz2020}. Higher pH and STP raised DRP losses, and greater
clay content lowered subsurface DRP losses, in a US--Canada synthesis; rainfall sharply
increased subsurface DRP loss where STP exceeded 41~mg~kg$^{-1}$
\parencite{Ni2020}. A change point in STP marks a threshold above which
CaCl$_2$-extractable P---and DRP in runoff and drainage---rises more steeply per unit
STP \parencite{McDowell2001b}. Drainage DRP was of similar magnitude to CaCl$_2$-P of
lysimeter topsoils, and water and CaCl$_2$ extractions were proposed to estimate runoff
and drainage DRP, respectively \parencite{McDowell2001b}. DRP concentration related to
soil P-AL in Swedish fields \parencite{Ulén2016}, high STP was an important spatial
driver of loss at the plot scale in Illinois \parencite{deOliveira2022}, and STP
strongly governed P desorption risk in flow-through experiments
\parencite{Mumbi2024}. Agronomic soil-P tests (Modified Morgan's P) related
significantly to drainage TDP and DUP across sites, with a steeper DRP relationship for
low-organic-matter cropland; batch extractions released far more soluble P than field
drain waters \parencite{Stutter2018}.

\subsubsection{Sorption capacity and degree of P saturation}

Where soils sorb P weakly, dissolved losses rise. Soils grouped as lower-P-sorbing had
distinctly higher tile-drainage TP and TPP (0.06 and 0.04 vs.\ 0.03 and 0.02~mg~L$^{-1}$)
and consistently exceeded the 0.03~mg~L$^{-1}$ standard; A-horizon degree of soil P
saturation (DSPS$_{M3}$) alone explained 83\% of TP variance, rising to 95\% with
water-extractable P, and grouping by sorption capacity increased predictive power to
96\% \parencite{Beauchemin2003}. In sandy German soils, rapid adsorption controlled
median drainage SRP: equilibrium concentrations and a DPS threshold of 0.25 were good
predictors, DPS exceeded 0.25 in all topsoils, and SRP $>$100~$\mu$g~L$^{-1}$ was
detected to great depth, showing local P transport and loss \parencite{Siemens2004}.
Across Norwegian profiles just 10~km apart, DPS-AL ranged 1--26\% and DPS-OX 3--37\%; a
whole-profile DPS-AL threshold near 10\% was suggested for TDP leaching, flow-weighted
TDP and DRP were best predicted by topsoil water-soluble P and whole-profile DPS-AL,
and the highest-leaching profiles combined large labile topsoil P pools with
low-sorption subsoils---evidence that subsoil characteristics must be included in
leaching risk assessment \parencite{Liu2024b}. Consistent with weak-sorption risk,
subsurface P losses were often small in fine-textured soils at moderate STP but became
environmentally significant in sandy soils, organic soils, and very high-P soils from
long-term fertilisation or manuring \parencite{Sims1998}. Adding alum to a
poor-P-sorbing soil reduced subsurface filtered reactive and total filtered P losses by
26--27\% at 50~kg~Al~ha$^{-1}$ (estimated cost 190--952~USD~kg$^{-1}$~P), with
mitigation rising with rate but declining over time \parencite{McDowell2015}. Anion
storage capacity was proposed to flag marginal soils with high P-loss potential
\parencite{McDowell2015b}.

\subsection{Soil Phosphorus Pools, Speciation, and Legacy Phosphorus}

Sequential and chemical fractionation link soil P pools to loss. P bioavailability and
loss are controlled by reactions with hydrous oxides and amorphous/crystalline Al, Fe,
and Ca complexes plus organic matter \parencite{Sharpley1995}. Sequentially extracted P
fractions correlated with P in overland flow and drainage, most strongly for anion
exchange membrane, bicarbonate, and HCl-extractable fractions, and isotopically
exchangeable P within 1~min correlated most closely; P lost to water was associated with
short-term (\textless24~hr) exchangeable pools linked to soil Al, Fe, and Ca, whereas P
exchangeable only after long periods was unrelated to solution P
\parencite{McDowell2001}. In organic soils, Ca-bound P was the largest pool and acted as a
P sink; under controlled drainage TP correlated most with HCl-P ($r=0.76$) and DRP with
water-extractable P ($r=0.81$), whereas under pipe drainage both correlated most with
NaOH-P \parencite{Grenon2022,Grenon2021}. Depressional (closed-depression) soils in
tile-drained watersheds were P-cycling hotspots: water-extractable, Mehlich-3, and total P
and all Hedley fractions were greater in depressions than hillslopes, decreasing with
depth, with NaOH-P and HCl-P the main pools; management history (P application and tillage)
governed concentration, DPS, and stratification \parencite{Mumbi2024}. Cultivated peatland
had the highest P and P-binding element contents, dominated by inorganic P (NaOH-Pi in plow
layers, H$_2$SO$_4$-Pi in subsoil), whereas abandoned, forest, and pristine peat retained
more organic and residual P \parencite{Höyhtyä2025}. Drainage sediments themselves store
substantial P---predominantly non-occluded Al/Fe-bound P---with high retention capacity
correlated to oxalate-extractable Fe, total C, and pH, and can release P under some
conditions even when water concentrations are low \parencite{Nguyen2002}. Legacy P is
manifest at watershed scale: cumulative SRP load gradients in Lake Erie tributaries rose
94--253\% after 2002 while cropland P balances were drawn down by 17--39~kg~P~ha$^{-1}$ over
15 years, with 63--66\% of the post-2002 SRP increase attributable to greater delivery
rather than higher runoff, consistent with increased tile drainage and connectivity
\parencite{Jarvie2017}.

\subsection{Organic Soils, Peatlands, and Redox-Mediated P Dynamics}

Organic and drained peat soils are distinctive high-risk P sources. They have been linked
to excessive P leaching and large TP loads, with water movement---modulated by
wetting/drying, organic-matter degradation, hydrophobicity, and shifting flow paths---the
key transport driver \parencite{Grenon2021}. Flooding organic soils increased P release to
drainage by 4--8 times relative to drained soils, and organic-P mineralisation can supply
up to 25~kg~P~ha$^{-1}$~yr$^{-1}$ in temperate mineral soils and up to
160~kg~P~ha$^{-1}$~yr$^{-1}$ in organic soils \parencite{Sharpley1995}. Very large filterable
P loads were measured from an Organic soil (86.7~kg~P~ha$^{-1}$) versus intergrade (9.0) and
Podzol (1.7), with added P leaching throughout the Organic profile but stratified in the top
15~cm of the Podzol \parencite{McDowell2015b}. In Finnish cultivated peat, thicker peat
increased dissolved-P share and load, dissolved P was higher in thick than thin peat plots
(38 vs.\ 22~$\mu$g~L$^{-1}$), and mineral subsoil retained P better than peat subsoil, so
drains set in mineral subsoil may capture P more effectively than those under thick peat
\parencite{YliHalla2022,Pham2023,Höyhtyä2025}. Degraded, artificially drained peatlands were
identified as the highest-risk eutrophication source in a German catchment, where grassland
peat P dynamics were governed by seasonal, redox-mediated remobilisation rather than
discharge alone, and fast flow exported most TP at tile-drained sites
\parencite{Tiemeyer2009}. Redox and iron chemistry directly mediate P: storm-triggered
mycelial ochre flocs from tile drains carried 27--141~mg~L$^{-1}$ TP with Fe and P
co-distributed along fungal mycelia and active Fe oxides forming 59--94\% of total Fe,
implying large episodic P discharge \parencite{Zimmer2016}. In boreal acid sulfate soils,
the redox interface lay at a median 1.6~m, with soil texture exerting stronger control than
drainage method (faster oxidation in sandy/silty than clayey soils)
\parencite{Virtanen2017}; reactive-transport modelling reproduced the geochemistry only when
both pyrite and faster-oxidising mackinawite and cation exchange were included, and increased
drainage efficiency and steady flow promoted sulfide oxidation above drain depth
\parencite{Salo2023}.

\subsection{Temporal Dynamics: Seasons, Events, and Concentration--Discharge Behaviour}

Phosphorus loss is highly concentrated in time. P concentrations and losses in tile drainage
are highly variable in space and time; losses are typically greater during storms than
baseflow, and most annual loss occurs in the non-growing season
\parencite{King2015}. In Ontario, the non-growing season accounted for 83--97\% of runoff,
84--100\% of DRP, and 67--98\% of TP loss \parencite{VanEsbroeck2016}, and models and field
data likewise placed most drainage volume and TP export in the non-growing season
\parencite{Qi2018,Sadhukhan2019b}. Event dominance is stark: more than 75\% of DRP and 66\% of
TP export occurred in the top 25\% of flows, and a single 6-day event supplied 46.6\% of TP,
84.0\% of DRP, and 81.0\% of suspended solids in one subwatershed \parencite{Brendel2019}; a
cited study found 88\% of TP loss in only 10\% of observation time \parencite{Elsayed2023};
and flooding/high-flow events dominated tile P loss in Illinois \parencite{Li2020b}.
Concentration--discharge behaviour differs by pathway: tile TP and DRP increased with
discharge (chemodynamic) while surface runoff was more chemostatic, and the tile DRP:TP ratio
fell at high flow, indicating a growing PP contribution \parencite{Osterholz2020}. Dissolved
silica in tiles was chemostatic but transport-limited at watershed outlets
\parencite{Sethna2022}. Seasonally, soluble P peaked in growing months in tile-intensive
Illinois \parencite{Algoazany2007}, whereas snowmelt and winter thaws dominated exports in
cold regions, where nutrients are exported predominantly in dissolved form---limiting the
value of particulate-focused practices \parencite{Baulch2019,VanEsbroeck2016}. Loads were
highest during and after snowmelt in Michigan \parencite{Clement2017}, and extreme events
after fertiliser application drove episodic export in Massachusetts and Indiana, where a
drought-then-deluge sequence in 2021 produced unprecedented P export (43\% of a 12-year DRP
load in nine months) \parencite{Millar2025,Mumbi2025b}.

\subsection{Management and Fertilisation History Controls}

\subsubsection{Fertilisation and manure history}

Long-term P inputs raise subsurface P losses. Long-term fertilisation increased DRP, dissolved
unreactive P, PP, and TP in tile drainage, generally in the order continuous grass $>$ rotation
$>$ continuous corn, and fertilised continuous grass shifted from a P sink to a significant P
source---implying that P removal by harvest is needed to limit losses
\parencite{Zhang2015b}. Swine-manure compost produced much greater DRP, PP, and TP than control
or yard-waste compost regardless of water-table management, with 4-year cumulative DRP up to
25.2~kg~ha$^{-1}$ and annual TP up to 8.36~kg~ha$^{-1}$~yr$^{-1}$ \parencite{Zhang2015}. Manure
source and form matter: solid cattle manure reduced flow, PP, and TP losses relative to
inorganic P and liquid manure despite similar DRP loads, and liquid manure gave the highest
immediate post-application DRP \parencite{Wang2018}. Organic amendments raised DRP loss more
under no-till than conventional tillage (a 113\% larger increase), driven solely by
concentration under no-till but by concentration and flow under conventional tillage
\parencite{Zhang2017}. Losses generally rose where soils were P-enriched, application rates
exceeded recommendations, or applications preceded large storms; adhering to soil-test-based
rates lowered loss, though organic vs.\ inorganic source and broadcast vs.\ subsurface placement
did not differ significantly in some Ohio comparisons \parencite{King2018}. Incidental losses
after surface application are important---about 60\% of DRP losses from a no-till cover-crop site
followed fall surface fertilisation \parencite{Macrae2023}, rainfall shortly after application
caused sharp Dutch pipe-drain P spikes \parencite{vanderSalm2012}, and a single manure event
supplied 8\% of a 12-year DRP load in Indiana \parencite{Mumbi2025b}. Across the MANAGE database,
generally less than 2\% of applied P was lost in drainage, loss rose only weakly with rate, and
organic vs.\ inorganic applications did not differ significantly \parencite{Christianson2016}.

\subsubsection{Tillage, cover crops, and rotation}

Tillage and residue management involve trade-offs between particulate and dissolved P.
Conservation tillage, no-till, and cover crops were associated with higher subsurface DRP
losses, while tillage increased PP and TSS \parencite{Ni2020,Macrae2023}. No-till DRP losses were
consistently greater and $\sim$60\% incidental, and lower PP at the no-till cover-crop site only
partly offset higher DRP so that TP was only slightly lower there \parencite{Macrae2023}. Reduced
tillage cut erosion but markedly increased DRP, so it should be paired with efficient subsurface
drainage to minimise runoff \parencite{Koskiaho2002}; no-tillage had significantly greater
dissolved P drainage loads than conventional tillage in the MANAGE synthesis
\parencite{Christianson2016}. Cover-crop effects on P are inconsistent: a cereal rye cover crop
did not significantly reduce drainage discharge or DRP in Illinois \parencite{deOliveira2022},
cover crops tended to raise DRP and TP leaching in lysimeters but not significantly at field drain
depth \parencite{Liu2019}, yet modelled winter rye reduced DRP by 16\% (mostly in runoff)
\parencite{Grewal2025} and an oat cover crop reduced predicted DRP and TP by 30\% and 54\%
\parencite{Askar2021b}. Corn--soybean and corn--soybean--wheat rotations reduced subsurface DRP
relative to continuous corn \parencite{Ni2020}, and DRP losses were greater in corn years than
soybean years \parencite{Zhang2015}. Simulations indicated moldboard/conventional tillage could reduce
drainage DRP and PP (by up to 48\% and 30\% as tillage intensity rose from no-till to
moldboard plow) while simultaneously lowering runoff P losses: under one compost scenario,
conventional tillage decreased overall runoff P loss by $\sim$50\% and tile-drainage P loss by
$\sim$19\% relative to no-till \parencite{Pan2023}, and shifting tillage to spring or injecting fertiliser could cut TP loss by
$\sim$10\% and 25\% \parencite{Qi2018}. In snowmelt regions, however, zero tillage and perennial
forages can \emph{increase} dissolved P despite reducing erosion---alfalfa forage lost 160\% more
TP than cropland, largely via a 221\% rise in dissolved P---so late-fall/winter applications are
consistently high-risk \parencite{Baulch2019}.

\subsubsection{Drainage design and grazing}

Drainage geometry alters both amount and pathway of loss. Increasing drain spacing reduced
subsurface volumes and TP loss (more strongly on finer soils) while increasing surface runoff and
its TP load \parencite{Morrison2013}. Shallower, closer drains increased P-loss risk in paddy
systems \parencite{DarziNaftchali2016}, and shallow tiles and surface inlets are among the field
attributes prone to elevated tile P \parencite{King2015}. Grazing intensity sharply raised losses:
a single intensive dairy grazing event increased sediment 50\%, particulate-associated P 40-fold,
and dissolved P 15-fold in tile water \parencite{Gillingham2000}, and grass--clover pasture leached
more subsurface MRP and TP than monoculture grass \parencite{Hooda1999}. Poorly drained,
high-STP sites had large potential to lose DRP in overland flow, with an 11-fold higher
flow-weighted DRP concentration from a high- versus low-fertility grassland
\parencite{Kurz2005}.

\subsection{Phosphorus Indices and Process Models for Drainage P}

Phosphorus indices are semiquantitative tools that delineate critical source areas by combining
source and transport factors, and their scores generally correlate positively with measured or
simulated losses across scales, though many region-specific versions exist and the more
physically based ones (e.g.\ the Swedish index) demand more data
\parencite{Buczko2007}. The Norwegian P index explained 66\% of variation in relative TP
concentration and identified the highest-risk fields and subcatchments, with soil P status
explaining 66\% of the source factor and erosion risk and contributing distance driving transport
\parencite{Bechmann2007}. A recurring critique is that most indices ignore tile drains or treat
them crudely; component-index frameworks have been proposed to represent tile P export explicitly
through source\,$\times$\,transport terms, accounting for the macropore-flow fraction of
precipitation and the dissolved and particulate concentrations it carries, with MF identified as
the dominant uncertainty and seasonal export still to be incorporated
\parencite{Reid2012}. Catchment nutrient-loss indices commonly use contributing distance,
connectivity, soil properties, and erosion; P is chiefly controlled by overland-flow conduits and
barriers whereas subsurface flow controls N more, and incorporating riparian removal capacity is
recommended \parencite{Pärn2012}. A simplified P index ranked subsurface treatments as lower risk
and surface-drained paddy as highest \parencite{DarziNaftchali2016}.

Process models remain limited for drainage P but are advancing. Few reviewed models address P in
drainage: SWAT includes tile flow but not tile P, DRAINMOD simulates drains but not P, and
ICECREAMDB---designed specifically for drainage P---is promising but needs testing; nested,
paired field experiments and uncertainty assessment are called for
\parencite{Radcliffe2015,Wang2020b}. DRAINMOD-P, developed to represent P dynamics in drained
cropland, satisfactorily predicted subsurface discharge (monthly NSE 0.59) and TP but poorly
predicted monthly DRP, attributing $\sim$21\% of tile DRP to macropore flow ($\sim$15\% of
discharge) \parencite{Askar2021,Askar2021b}. ICECREAM satisfactorily simulated hydrology, TP, and
DRP but fared worse for tile PP and soil moisture in Ontario, where observed tile TP was 65\% of
total \parencite{Qi2018}. RZWQM2-P---using Richards, Hooghoudt, and Poiseuille formulations for
matrix, tile, and macropore flow---simulated surface and tile DRP and PP and total P well, was the
first to accurately simulate daily DRP, and reproduced pathway partitioning
\parencite{Sadhukhan2019,Sadhukhan2019b,Grewal2025}. EPIC reasonably modelled flows and tile DRP,
and coupling with SurPhos improved runoff DRP and manure representation
\parencite{Wang2019}. RZWQM2-P also captured long-term tillage effects on tile P
\parencite{Pan2023}, and DRAINMOD linked to empirical regressions predicted tile P loads with
minimal parameters and confirmed that wider drain spacing reduces subsurface P loss
\parencite{Morrison2013}. Models such as SurPhos and manure-runoff and direct-fertilizer models
have shown accuracy for runoff P, while ICECREAM performed poorly for soil moisture and tile PP
\parencite{Wang2020b}.

\subsection{Edge-of-Field Retention and Mitigation}

The soil and pathway controls above define where mitigation must act, and a consolidated body of
evidence quantifies edge-of-field practices---while emphasising trade-offs and site specificity.
Drainage water (controlled) management can reduce loss by curtailing discharge: it cut dissolved P
by 63\% and TP by 50\% in Minnesota \parencite{Kröger2013}, and managed drainage reduced tile
ortho-P loss by 80\%, with more than half the reduction occurring in spring at similar drainage,
suggesting enhanced crop P uptake \parencite{Nash2015c}. Controlled drainage with subirrigation
reduced DRP, PP, and TP for low-P treatments but not for swine-manure compost
\parencite{Zhang2015}, and reviews report TP reductions of 57\% and soluble-P reductions of
57--80\% \parencite{Elsayed2023}; however, controlled drainage can also raise DRP/TP or shift loss
to surface runoff, and organic-soil concentrations may stay above limits even as loads fall
\parencite{Grewal2025,Sadhukhan2019b,Grenon2021}. Blind inlets replacing tile risers reduced TP by
up to 78\% and soluble P by up to 72\% \parencite{Kröger2013}, with 50--78\% TP and 65--72\% DRP
reductions reported elsewhere \parencite{Singh2020c}. Subsurface placement and 4R management are
highlighted as promising, with deep placement cutting DRP by 69--72\%
\parencite{Wilson2019,Sharpley2019,Singh2020c}. Denitrifying bioreactors gave modest, variable DRP
removal (median 12\% efficiency; outflow DRP less than half inflow) and occasionally acted as DRP
sources under low summer flow \parencite{Bailon2025}. Filter media (Fe- and Al-based superior to
Ca/Mg) and passive Fe-oxide resins target dissolved P, and P-sorbents can extend system lifetime
under oxidising conditions \parencite{Mendes2020,Li2020b}. Constructed and vegetated systems retain
P chiefly by particulate sedimentation, with retention of 27--100\% but strong dependence on
hydraulic residence time, sorbent availability, and loading---systems can flip from sink to source
during events \parencite{Mendes2020,Kröger2013,Singh2020c}. Vegetated ditches retained
$\sim$44\% of inorganic P annually \parencite{Kröger2013}, drainage sediments and ditch beds shifted
seasonally between P sink, source, and equilibrium \parencite{Ahiablame2010,Nguyen2002}, and routing
$\sim$25\% of tile flow through a riparian buffer achieved near-complete removal of dissolved
nutrients and sediment \parencite{Jacquemin2020}. Because buffers can bypass the root zone via tile
loads and may saturate over time, modified, bypass-preventing designs and trap-and-treat ditches
are emphasised \parencite{Mendes2020,Millar2025,Baulch2019}. Finally, tile drainage does not
inevitably raise P export: under some conditions installing tiles reduced surface runoff and SRP
losses without increasing TP export \parencite{Klaiber2020}, and in frozen vertisolic landscapes
tile expansion may not exacerbate P (though it can raise N) \parencite{Kokulan2021}. Overall,
effective P control depends on matching the right practice to site hydrology, connectivity, weather,
and legacy P, and on recognising N--P trade-offs among conservation practices
\parencite{King2015,Sharpley2019,Ni2020,Baulch2019,Wilson2019}.

\printbibliography
\end{document}
